% ============================================================
% PREAMBLE.TEX — Packages et configuration
% ============================================================

% ------------------------------------------------------------
% ENCODAGE ET LANGUE
% ------------------------------------------------------------
\usepackage[utf8]{inputenc}
\usepackage[T1]{fontenc}
\usepackage[french]{babel}
\usepackage{lmodern}
\usepackage{microtype}

% ------------------------------------------------------------
% MISE EN PAGE
% ------------------------------------------------------------
\usepackage{geometry}
\geometry{
    margin=2.5cm,
    headheight=15pt
}

% ------------------------------------------------------------
% GRAPHIQUES ET COULEURS
% ------------------------------------------------------------
\usepackage{graphicx}
\graphicspath{{images/}}
\usepackage{xcolor}
\usepackage{tikz}
\usepackage[most]{tcolorbox}

% ------------------------------------------------------------
% TABLEAUX
% ------------------------------------------------------------
\usepackage{booktabs}
\usepackage{array}
\usepackage{multirow}
\usepackage{longtable}

% ------------------------------------------------------------
% LISTES
% ------------------------------------------------------------
\usepackage{enumitem}
\setlist[itemize]{itemsep=2pt,parsep=0pt}
\setlist[enumerate]{itemsep=2pt,parsep=0pt}

% ------------------------------------------------------------
% TITRES DE CHAPITRES ET SECTIONS
% ------------------------------------------------------------
\usepackage{titlesec}

% CORRECTION 1 : utiliser \@chapapp (qui vaut "Chapitre"
% puis "Annexe" après \appendix) au lieu de la chaîne
% "Chapitre" codée en dur. \makeatletter est nécessaire
% pour tokeniser \@chapapp correctement.
\makeatletter

\titleformat{\chapter}[display]
  {\normalfont\huge\bfseries}
  {\filright\Large \@chapapp\ \thechapter}
  {1ex}
  {\titlerule\vspace{1ex}\filleft}
  [\vspace{1ex}\titlerule]

% CORRECTION 2 : variante "numberless" pour \chapter*,
% sinon titlesec imprime "Chapitre 0" sur l'Avant-propos,
% la Table des matières, la bibliographie non numérotée, etc.
\titleformat{name=\chapter,numberless}[display]
  {\normalfont\huge\bfseries}
  {}
  {0pt}
  {\titlerule\vspace{1ex}\filleft}
  [\vspace{1ex}\titlerule]

\makeatother

\titleformat{\section}[hang]
  {\normalfont\Large\bfseries}
  {\thesection.}{1em}{}

\titleformat{\subsection}[hang]
  {\normalfont\large\bfseries}
  {\thesubsection}{1em}{}

\titleformat{\subsubsection}[hang]
  {\normalfont\normalsize\bfseries}
  {\thesubsubsection}{1em}{}

% ------------------------------------------------------------
% EN-TÊTES ET PIEDS DE PAGE
% ------------------------------------------------------------
% CORRECTION 3 : \nouppercase est fourni par le paquet
% "textcase", PAS par fancyhdr ni par la classe book.
% Sans ce paquet, la commande est indéfinie → erreur.
\usepackage{textcase}

\usepackage{fancyhdr}
\pagestyle{fancy}
\fancyhf{}
\fancyhead[L]{\small\nouppercase{\leftmark}}
\fancyhead[R]{\small\thepage}
\renewcommand{\headrulewidth}{0.4pt}

% ------------------------------------------------------------
% MATHÉMATIQUES (si nécessaire)
% ------------------------------------------------------------
% CORRECTION 4 : amsmath DOIT être chargé AVANT hyperref
% (hyperref patche les environnements d'équations d'amsmath ;
%  chargé après, \eqref et \tag ne créent pas de liens corrects).
\usepackage{amsmath}
\usepackage{amssymb}

% ------------------------------------------------------------
% HYPERLIENS — à charger EN DERNIER (avant les commandes perso)
% ------------------------------------------------------------
\usepackage{hyperref}
\hypersetup{
    colorlinks=true,
    linkcolor=black,
    urlcolor=blue,
    citecolor=black,
    pdftitle={Cours de design graphique},
    pdfauthor={Ton nom},
    bookmarksnumbered=true
}

% ------------------------------------------------------------
% COMMANDES PERSONNALISÉES
% ------------------------------------------------------------
\newcommand{\terme}[1]{\textbf{#1}}
\newcommand{\important}[1]{\textbf{\textcolor{red!70!black}{#1}}}
\newcommand{\logiciel}[1]{\textsf{#1}}
\newcommand{\anglais}[1]{\textit{#1}}

% Encadrés pédagogiques
\newtcolorbox{retenir}{enhanced,breakable,colback=blue!4,colframe=blue!60!black,
  fonttitle=\bfseries,title={À retenir},left=6pt,right=6pt}
\newtcolorbox{astuce}{enhanced,breakable,colback=green!5,colframe=green!45!black,
  fonttitle=\bfseries,title={Astuce de pro},left=6pt,right=6pt}
\newtcolorbox{erreur}{enhanced,breakable,colback=red!4,colframe=red!60!black,
  fonttitle=\bfseries,title={Erreur fréquente},left=6pt,right=6pt}
\newtcolorbox{defbox}[1][Définition]{enhanced,breakable,colback=orange!5,colframe=orange!70!black,
  fonttitle=\bfseries,title={#1},left=6pt,right=6pt}